% id: 132-01
% book: 베이직쎈 확률과 통계 · 07 통계적 추정
% section: 개념 40 모평균의 추정
% figure: none
% answer: 
% answer_source: 
% variant_level: 0 (원본 전사)
% 이 파일은 build.mjs 가 items.json 에서 만든다. 고칠 때는 items.json 을 고친다.
\smash{\raisebox{1.5mm}{\dmkichul{베이직쎈 확률과 통계 132쪽 01번}}}
다음은 정규분포 $\mathrm{N}(m{,}\mkern3mu\mathopen{}\sigma^2)$을 따르는 모집단에서 크기가 $n$인 표본을 임의추출할 때, 모평균 $m$에 대한 신뢰도 $95\mkern3mu \%$의 신뢰구간을 구하는 과정이다. ㈎, ㈏, ㈐에 알맞은 것을 구하시오. \exprbox{\begin{minipage}{0.96\linewidth}\raggedright 표본평균 $\overline{X}$는 정규분포 $\mathrm{N}\!\left(m{,}\mkern3mu\mathopen{}\fbox{㈎}\right)$을 따르므로 $Z=\dfrac{\overline{X}-m}{\fbox{㈏}}$으로 놓으면 확률변수 $Z$는 표준정규분포 $\mathrm{N}(0{,}\mkern3mu\mathopen{}1)$을 따른다.\\ 이때 $\mathrm{P}(0\le Z\le 1.96)=0.475$에서\\ \hspace*{1em}$\mathrm{P}(-1.96\le Z\le 1.96)=2\times 0.475=0.95$\\ 즉 $\mathrm{P}\!\left(-1.96\le\dfrac{\overline{X}-m}{\fbox{㈏}}\le 1.96\right)=0.95$이므로\\ \hspace*{1em}$\mathrm{P}(\overline{X}-\fbox{㈐}\le m\le\overline{X}+\fbox{㈐})=0.95$\\ 따라서 표본평균 $\overline{X}$의 값이 $\overline{x}$일 때, 모평균 $m$에 대한 신뢰도 $95\mkern3mu \%$의 신뢰구간은\\ \hspace*{1em}$\overline{x}-\fbox{㈐}\le m\le\overline{x}+\fbox{㈐}$\end{minipage}}
